\chapter{Zusammenfassung und Ausblick} \label{chpt:Zusammenfassung_und_Ausblick_Main}

Am Ende der Arbeit wird das erreichte nochmals zusammengefasst um von diesem Standpunkt aus einen Blick in die Zukunft zu wagen.

Rückblickend lässt sich sagen, welche Entscheidungen zielführend und welche Entscheidungen hinderlich waren. Dies wird in der Zusammenfassung diskutiert.

Für die Zukunft werden im Ausblick Ideen angeführt, die während oder nach der Arbeit an der Lösung entsanden sind. Diese Gedanken sollen als Wegweiser für eine eventuelle Verbesserung oder Weiterführung der Arbeit dienen.




\section{Zusammenfassung} \label{sec:summary}

Im Rahmen dieser Arbeit wurde die existierende NEORV32 CPU um Komponenten zur Matrix-Berechnung erweitert. Die Komponenten unterliegen starken Einschränkungen. Unter anderem können zum jetzigen Zeitpunkt lediglich 4x4 Matrizen mit vorzeichenlosen 32-Bit Zahlen berechnet werden. Zum termingerechten Abschluss dieser Arbeit war es notwendig solche Einschränkungen vorzunehmen. Für die Zukunft wäre es durch zusätzlichen Arbeitseinsatz  möglich das Design um weitere Datentypen und Konvertierungen zu erweitern.

Die durchgeführten Änderungen orientieren sich an der Befehlsmenge der bereits existierenden RISC-V Vektor-Extension und greifen damit auch das Konzept des Strip-Mining auf. Strip-Mining erlaubt die Erstellung eines Hardware-Software-Co-Designs, welches in der Lage ist Matrixmultiplikationen größere Matrizen durch Multiplikationen kleinerer Dimensionen zusammenzusetzen.

Um die Softwareentwicklung für Anwendungen und Testprogramme zu ermöglichen, wurden exzessive Anstrengungen unternommen, die etwa den gleichen Zeitaufwand wie die VHDL Entwicklung zur Folge hatten. Insbesondere wurde es durch einen Compiler ermöglicht die neuen Anweisungen auch unter Verwendung einer höheren Programmiersprache wie C in RISC-V Maschinencode zu übersetzen. Im Nachhinein wäre der schnellstmögliche Einsatz der NEORV32 Intrinsics ein sehr großer Vorteil gewesen. Leider wurde dieser Weg erst spät gelernt.

Es wurde ein Vergleich zweier Testprogramme auf dem NEORV32 Chip durchgeführt. Die Messungen zeigen einen Geschwindigkeitszuwachs um einen relevanten Faktor. Das Experiment ist geglückt und könnte in Zukunft durch die Erweiterung auf ein vollständiges Typsystem und weiter verbesserte Speicherzugriffe noch verbessert werden.



\section{Ausblick} \label{sec:outlook}

Am Ende dieser Arbeit drängt sich eine grundlegende Frage auf. In welchen Szenarion ist es sinnvoll Hardwarebeschleuniger in ein bestehendes RISC-V Design einzupflegen und ab wann sollte sich ein Unternehmen für die Erstellung eines komplett neuen Designs entscheiden. Diese Frage ist eng verknüpft mit der Frage ob der in dieser Arbeit gewählte Ansatz richtig oder falsch ist.

Im Zuge dieser Arbeit wurde ein bestehendes Design erweitert. Der Zugriffsweg auf den Speicher über die Load-Store-Unit und den am Bus angeschlossenen Speicher wurde wiederverwendet. Der Vorteil ist es sicherlich Kosten und Zeit sparen zu können, da existiernde Komponenten wiederverwendet werden. Ein Nachteil ist eventuell, dass die wiederverwendeten Komponenten auf Anforderungen hin erstellt und optimiert worden sind, die in der Vergangenheit maßgebend waren aber eventuell neue Erweiterungen nicht effizient unterstützen. Für diese Arbeit ist dieser Umstand nicht zu einem Problem geworden.

Die Alternative wäre es, ein grundlegend neues Design zu beginnen welches dann die neuen Anforderungen besser abdeckt. Ein Beispiel für einen neue Anforderung ist es beispielsweise mehr als einen Speicher-Zugriff parallel ausführen zu können. Die Hardware für die Outer-Produkt Berechnung könnte dann mehrmals synthetisiert werden. Jede Instanz der Outer-Produkt Berechnung könnte möglicherweise über ihren eigenen Zugriff auf den Speicher verfügen. Wenn der Speicher echt-parallel Daten über mehrere Ports gleichzeitig liefern kann, wird eine Matrix-Berechnung parallel statt sequentiel ausgeführt.

Ein Nachteil eines solch spezialisierten Designs ist es, dass es seine Allgemeinheit verliert. Ein RISC-V CPU soll eine allgemeine Rechenmaschine bleiben.

Es ist fallabhängig, welchen Weg ein Unternehmen wählen sollte. Ein finanzstarkes Unternehmen kann sich für die Erstellung von Anwendungsspezifischen Chips entscheiden, wie man es im Falle von Amazon und dem Trainium Chip sehen kann. Unternehmen wie NVIDIA erstellen Beschleunigerkarten, die für High-Performance Computing im Sinne von großer Parallelität eingesetzt werden können und damit ein großeres Feld abdecken. CPU Designs von Intel und AMD sind um Instruktionserweiterungen wie AVX2, AMX und FMA für die Vektor- und Matrixberechnung erweitert worden wobei sie immer noch allgemeines Computing erlauben. In Zukunft wird RISC-V diesen Weg mit einer Matrix-Extension auch gehen und hat mit der \gls{RVV} bereits begonnen.

%\footnote{https://aws.amazon.com/de/ai/machine-learning/trainium/}

Wahrscheinlich ist es aus wirtschaftlicher Sicht notwendig zu wissen, was der Markt fordert und Hardware erstellen zu können die die Anforderung des Marktes erfüllen. Aus akademischer Sicht ist es wichtig die Zusammenarbeit mit Firmen aus der Industrie zu suchen um innovative Lösungen für die zukünftigen Entwicklungen finden zu können.

Als nächste Schritte zur Fortführung dieser Arbeit ist es denkbar die Matrix-Extension intelligenter und damit performanter zu gestalten. Der Vergleich mit anderen Designs wie z.B. \cite{DBLP:journals/corr/abs-2104-03142} und \cite{scoss} führt zu neuen Erkenntnissen und einen Zuwachs an Erfahrung.

Zusätzlich ist das System aus dieser Arbeit um weitere Datentypen zu ergänzen. Es gibt traditionelle Datentype wie beispielsweise 8-, 16-, 32- und 64-bit ganzzahlige Datentypen mit und ohne Vorzeichen und auch 32- und 64-bit Gleitkommadatentypen. Dazu kommen neuartige Aufteilungen von Daten, wie sie für Large Language Models verwendet werden. Ein Beispiel ist hier NF4, was für normalized float mit 4-Bit steht. Sobald die Matrix-Extension mit unterschiedlichen Datentypen verwendet werden können soll ist über die Konvertierung zwischen den Datentypen nachzudenken. Es müssen ebenfalls dazu passende Instruktionen ergänzt werden. Die set(i)vli Instruktion aus der \gls{RVV}-Extension benötigt dann ein Äquivalent in der Matrix-Extension um die Matrix-Extension für Datentypen konfigurierbar zu machen.

Es ist zu prüfen, welche Ergebnisse das Lowering hervorgebracht hat. Wurden wirklich Hardware-Bausteine innerhalb des FPGA verwendet oder sind Anteile unnötigerweise durch Logikzellen synthetisiert worden? Hier lässt sich in Kleinstarbeit noch eine Optimierung durchführen, falls Logikzellen anstatt von vorhandener Hardware verwendet werden.

Ein weiterer Schritt ist es sicherlich eine Umgebung für die Validierung des Systems bereitzustellen. Aufgrund zeitlicher Beschränkungen wurde das Design nur visuell mit einem Waveform-Viewer geprüft. Die Sprache System-Verilog beispeilsweise bietet weitrechendere Möglichkeiten zur Validierung.

%Für VHDL könnte die UVVM \footnote{https://www.uvvm.org/} oder cocotb \footnote{https://www.cocotb.org/} eingesetzt werden.